\chapter{The Idea: Numbers That Know How Much Room They Have}
\label{ch:idee}

% Level: surface. No displayed formulas, only numbers in running text.
% Respect the order of discovery: reversed bases -> finiteness ->
% leading zeros (correction of 13 July 2026).

\section{An ordering nobody questions}

Every positional number system rests on a silent convention. In the
decimal system the last digit counts ones, the one before it tens,
then hundreds -- each position weighs ten times its neighbour. The
binary system does the same in steps of two. And there is a less
familiar but venerable system in which the bases are not constant
but grow: the factorial system. There the last digit may only be 0
or 1, the next one 0 to 2, the one after that 0 to 3, and so on.
The further left a position sits, the more it may count \emph{and}
the more it is worth.

In all these systems the same direction points upward: the position
with the most freedom carries the most weight. This seems so
self-evident that nobody ever says it out loud.

Horimetrik begins with the question of what happens when you
reverse exactly that. What if, of all positions, the
\emph{tightest} one -- the one that knows only 0 and 1 -- gets the
greatest weight? What if the positions grow ever more generous
towards the right, while their influence keeps shrinking?

\section{Enforced finiteness}

The first consequence arrives immediately, and it is not a design
decision but a necessity. An ordinary number system grows leftward
into infinity: when a number gets too large, you prepend another
position with an even greater weight -- there is always room. In
the reversed system this is impossible. On the left already sits
the tightest position, the one with base two, and below two there
is nothing. The system is \emph{sealed} on the left.

Whoever works with a fixed number of positions therefore gets a
finite supply of numbers. With one position there are two, with two
positions six, with three twenty-four, with four one hundred and
twenty -- every additional position multiplies the supply, and the
sizes that arise are precisely the factorials. Once the supply is
exhausted, no carry will help: you need a \emph{new system} with
more positions, in which every old digit suddenly carries a
different weight.

This finite living space of a number is what we call its
\emph{horizon}. What at first looks like a defect of the reversed
system -- you cannot simply keep counting -- is in truth its most
interesting feature: every number here carries not only a value but
also the information of which space that value lives in. The
totality of all horizons remains infinite; only each single one is
finite. Infinity still exists -- but as a staircase, not as a
straight line.

\section{The second find: the zeros}

The leading zeros came later, and they came as a surprise. In every
familiar system a prepended zero is pure cosmetics: 012 is 12. In
the reversed system you prepend a zero -- and thereby change the
number.

An example with small numbers. In the horizon with two positions
the digit string 12 means the value five. Prepend a zero and it
reads 012 -- but now in a horizon with three positions, in which
all the weights are distributed differently. The value is six. One
more zero: 0012, value seven. The digits 1 and 2 have not moved,
and yet the number counts one further with every extension.

This is the moment at which a curiosity becomes an object of study.
For if a leading zero changes the value, then there are suddenly
\emph{two different answers} to the most harmless of all questions:
what does it mean to give a number more room?

\begin{itemize}
  \item You can hold its \emph{value} fixed -- then its digits must
        reorganize themselves when moving into the larger horizon.
  \item You can hold its \emph{shape} fixed -- the digits stay
        put, a zero is prepended -- then its value shifts.
\end{itemize}

In the decimal system these are the same operation. Here they are
two, and they lead into different worlds. Almost everything this
book studies lives in the gap between these two notions: the two
identities of a Hori number that are visible at first (a third will
reveal itself later), the ways of tidying up, the worlds of
arithmetic, and the tensions between them.

\section{What was actually invented here -- and what was not}

Honesty belongs here from the start: the building blocks of this
system are all known mathematics. Number systems with mixed bases
have existed for centuries, the factorial system is well studied,
and combinatorics has long known that its spaces have factorial
size. What is new -- at least we have found it nowhere in this form
-- is the combination: the reversed reading direction inside a
finite space, the horizon as part of the number, the
meaning-carrying zeros, and the two competing ways to grow. Whether
this combination can carry a theory of its own was the bet on which
this project began. The following chapters lay the cards on the
table.

\begin{oberflaeche}
If only one sentence of this chapter is to remain, let it be this:
Horimetrik did not arise from a wish to write numbers differently,
but from the question of what an ordering enforces once you reverse
it. The answer was: finiteness -- and from finiteness everything
else followed on its own.
\end{oberflaeche}
